\chapter{Pelindungan Data Pribadi dalam Praktik}

\section*{Tujuan Pembelajaran}
\begin{itemize}
  \item Menjelaskan kewajiban pengendali dan prosesor data pribadi menurut UU 27/2022 dan PP 33/2026.
  \item Menyusun catatan kegiatan pemrosesan data pribadi untuk satu proses layanan.
  \item Menganalisis insiden kebocoran data dan menilai apakah penanganannya memenuhi kewajiban notifikasi.
\end{itemize}

\section{Pendahuluan}

% TODO
% Cakupan: data pasien bocor lewat vendor, lalu muncul pertanyaan siapa yang bertanggung jawab dan apa yang wajib dilakukan dalam hitungan jam.
% Kasus Medikara: insiden kebocoran skala kecil lewat vendor pada 2025.

\section{Konsep Dasar UU 27/2022}

% TODO
% Cakupan: data pribadi umum dan data pribadi yang bersifat spesifik, termasuk data kesehatan; prinsip pemrosesan; dasar pemrosesan; ruang lingkup dan masa transisi yang berakhir 17 Oktober 2024.

\section{Hak Subjek Data dan Kewajiban Pengendali}

% TODO
% Cakupan: hak akses, pembetulan, penghapusan, dan keberatan atas keputusan otomatis; kewajiban pengendali dan prosesor; pengendali bersama; pejabat atau petugas pelindungan data pribadi.

\section{PP 33/2026 dan Status Penegakan}

% TODO
% Cakupan: ditetapkan 16 Juli 2026 dan berlaku 16 Januari 2027; rincian teknis seperti mekanisme persetujuan, data anak, penilaian dampak, transfer data, dan ganti rugi; sanksi administratif; status lembaga pengawas pelindungan data pribadi.
% Catatan: per 30 Oktober 2026 naskah resmi belum tampil di JDIH BPK; cocokkan nomor pasal dengan naskah resmi saat menulis isi.

\section{Penilaian Dampak dan Notifikasi Kegagalan}

% TODO
% Cakupan: kapan penilaian dampak pelindungan data pribadi wajib (data spesifik, skala besar, keputusan otomatis); isi minimum penilaian; notifikasi tertulis paling lambat 3 x 24 jam kepada subjek data dan lembaga; isi notifikasi.
% Kasus Medikara: garis waktu insiden vendor 2025 dibandingkan dengan kewajiban notifikasi.

\section{Catatan Kegiatan Pemrosesan dan Transfer Data}

% TODO
% Cakupan: struktur catatan kegiatan pemrosesan (RoPA, Records of Processing Activities); pemetaan alur data; transfer ke luar wilayah hukum Indonesia; kaitan dengan kerahasiaan dan masa simpan rekam medis menurut Permenkes 24/2022.
% Kasus Medikara: alur data pendaftaran pasien dan telekonsultasi di Medikara+.

\section{Sistem Manajemen Privasi ISO/IEC 27701:2025}

% TODO
% Cakupan: PIMS (Privacy Information Management System) yang kini dapat berdiri sendiri; kaitan dengan ISO/IEC 27001; kapan sertifikasi privasi berguna bagi organisasi.
% Pendalaman teknis: pseudonimisasi, tokenisasi, klasifikasi data otomatis, dan sumber daya Consent pada HL7 FHIR (dibahas rinci di IT40613).

\section{Aktivitas: Catatan Pemrosesan Data dan Analisis Insiden Vendor}

% TODO
% Tujuan aktivitas: memahami alur data pribadi yang sebenarnya dan menilai celah kepatuhan PDP berdasarkan satu proses dan satu insiden.
% Peran dalam tim: analis regulasi memimpin; pemilik konteks memetakan alur data nyata; analis teknologi menelusuri sistem yang terlibat; analis risiko memperbarui risk register.
% Langkah: pilih satu proses layanan; isi catatan pemrosesan (jenis data, tujuan, dasar, penerima, lokasi, masa simpan, kontrol); susun garis waktu insiden vendor 2025; nilai celah terhadap UU PDP dan PP 33/2026; rangkum seluruh luaran Bab 4 sampai 7 menjadi C2.
% Luaran dan tonggak capstone: catatan kegiatan pemrosesan data pribadi dan analisis celah PDP; masuk C2 yang ditinjau menjelang UTS.
% Berkas sources yang direncanakan: sources/chapter07/templat-ropa.xlsx, sources/chapter07/kronologi-insiden-vendor-2025.md, sources/chapter07/daftar-periksa-celah-pdp.xlsx

\section{Ringkasan}

% TODO

\section{Latihan}

\begin{enumerate}
  \item Medikara memakai vendor untuk aplikasi Medikara+. Dalam hubungan ini, siapa pengendali dan siapa prosesor? Apa akibatnya terhadap tanggung jawab saat terjadi kebocoran?
  \item Tentukan dasar pemrosesan yang paling tepat untuk data rekam medis dalam layanan pengobatan. Jelaskan mengapa persetujuan belum tentu menjadi dasar yang terbaik.
  \item Susun garis waktu notifikasi yang memenuhi UU PDP untuk insiden yang ditemukan pada Jumat malam.
\end{enumerate}
